\chapter{Open questions before one-shotting}
The specification checks, and its risky parts are checked by the server. What
follows is everything not yet in a Lean file: questions the source would
otherwise answer by guessing, requirements agreed but not yet written, and the
limits of the tools. Each item moves out of this chapter when it is realized.

The questions are grouped by what they hold up. The first group decides security
and must be settled before any source is written; the others shape features
that could ship later.

\section{Security: settle before the source}
\begin{enumerate}
\item \textbf{Keeping apps read-only.} An app runs in its viewer's browser, so
  whatever lets it read could let it write. The proposal: serve uploaded content
  from its own origin, whose session can only read, under a Content Security
  Policy that lets it reach nothing but this server; writes need credentials that
  origin never holds. Is that enough, or should each folder's apps get an origin
  of their own, so one app cannot read what another folder's app was shown?
\item \textbf{Who may write a policy?} Anyone who may write a folder can create
  \texttt{plans.doc--policy.rego} next to someone else's \texttt{plans.doc} and,
  being the last level of the walk, it overwrites what it likes. Creating or
  editing a \texttt{--policy.rego} should need \texttt{write} under the policy it
  would replace, and perhaps a separate decision: should \texttt{Eval} return
  \texttt{admin} beside \texttt{list}, \texttt{read} and \texttt{write}?
\item \textbf{May a lower level loosen?} \texttt{list} may not, because names are
  content. May \texttt{read}, below a folder the caller may list? May
  \texttt{write}? May a lower banner replace a higher one, or is the banner the
  greatest rank on the path?
\item \textbf{The root.} The walk starts at \texttt{/files} (and at
  \texttt{/apis}). What do their policies say, and who may write them? With no
  policy found at all, the answer is deny.
\item \textbf{A policy change reaches back in time.} A subscriber's cursor moves
  past records it may not read. If a later version of their policy lets the caller
  in, those records are already behind the cursor and are never sent; if it shuts
  the caller out, files already shown stay on screen. The new policy version is a
  record in the feed, but the subscriber may not be allowed to read it, so it
  cannot be the cue. The shard has to act: when a policy version commits, it
  evaluates the files below it again and sends the changed decisions as fresh
  events, appearing and withdrawing files. How is that bounded when the policy
  governs a large tree?
\item \textbf{Records relayed from an untrusted origin.} A peer's feed carries
  records it imported from third shards. If a shard trusts the peer but not the
  third shard, does it drop those records? A policy file from an origin it does
  not trust must not govern anything there. And so that a relay cannot forge an
  origin, should each record be signed by its origin shard?
\item \textbf{Attributes beyond the token.} APIs look users up by id, because not
  every permission is in the token. Where does that directory live, and what
  happens during a partition? If two shards hold different directory data, one
  token gets different answers on each, so the directory needs the same treatment
  as the records: replicated, and its answers as of the same records.
\item \textbf{Paths.} Is everything under \texttt{/files/} or \texttt{/apis/}? The
  upload endpoints (\texttt{POST /uploads}, \texttt{PUT /uploads/\{id\}}) sit
  outside both today; should an upload be a \texttt{PUT} to the file's own path
  under \texttt{/files/}, in parts?
  Names must not contain \texttt{/}, must not be \texttt{.} or \texttt{..}, and a
  name ending in \texttt{--policy.rego} is always a policy. The proposal for
  symlinks and relative paths, which seat folders tend to raise: none in the
  store; paths are canonical and \texttt{..} is refused, because a second name for
  a file would let a link grant access; ``my files'' across the seats a person
  fills is a view the UI builds from the seats in their token.
\item \textbf{Mounting.} A mount is an internal HTTP address that
  \texttt{ingressHttps} forwards to. How is it declared: a file in the tree, such
  as \texttt{/apis/db/.--mount.json}, written under the policy at that path, so
  mounts replicate like everything else? Does each island run its own copy of a
  service, or forward across islands? And how many requests may one mount have in
  flight, so that one slow service cannot take every connection the door has,
  the way a failing import could once take every lane?
\item \textbf{The cursor can reveal hidden activity.} When a page holds only
  records the caller may not read, the cursor still moves, which tells the caller
  that something they cannot see was written. Is that acceptable, or should the
  cursor advance only on a schedule, whatever was written?
\end{enumerate}

\section{Policy}
\begin{enumerate}
\item \textbf{What does \texttt{Eval} see?} Does each level's policy get the table
  from above as input, so it can decide relative to it (``as my parent, but not
  for contractors'')?
\item \textbf{Which entries are supplied?} Banner and colours are; what else
  (retention, content type, job list, allowed models)?
\item \textbf{Concurrent versions of one level's policy} are combined entry by
  entry: \texttt{list}, \texttt{read} and \texttt{write} by and, the banner by
  greatest rank. What about other entries?
\item \textbf{One rank for heterogeneous files.} A listing may show a medical
  record beside a national marking, and its page banner is the greatest rank among
  them. That needs every policy's rank on one scale. Who defines the scale: the
  root policy, or a supplied entry every policy must use? And what does the page
  banner say when the greatest rank belongs to a file of a different kind than
  the rest?
\item \textbf{Token age per tree.} Should a policy be able to supply a maximum
  token age, so a sensitive folder demands a token issued within the hour while
  the rest accept a day's?
\item \textbf{Tightening across a partition.} Should some classifications make
  policy changes wait for a quorum, so that a partitioned island cannot honour a
  policy that was tightened elsewhere?
\item \textbf{Break glass.} Medical systems let a clinician open a record they
  are not normally allowed to, in an emergency, on condition that it is recorded
  and reviewed. Is that a \texttt{read} granted with an audit record, and what
  does the banner say?
\item \textbf{Erasure.} GDPR gives a person the right to have their data erased,
  and this design never loses a version and replicates it everywhere. Erasing
  needs a tombstone that every shard honours by deleting the blobs of the named
  versions. Is that in scope?
\end{enumerate}

\section{Replication and storage}
\begin{enumerate}
\item \textbf{The same path written on both islands.} During a partition two
  people can upload \texttt{plans.doc} into one folder, one on each island. They
  get different \texttt{id}s, so after repair one path names two files. Is the
  later one renamed, are both shown, or does one become a version of the other?
  Folders, including a seat's home folder made the first time its token is seen,
  have the same problem, and one proposed answer: derive a folder's \texttt{id}
  from its path, so that \texttt{mkdir -p} on both islands makes the same folder,
  and two first policies become concurrent versions of one policy, which both
  must allow, and so fail safe.
\item \textbf{Clocks.} \texttt{dt} is nanoseconds at the origin, and the latest
  version is the greatest \texttt{(dt, shard)}, so an island whose clock runs fast
  wins every concurrent edit. Is that acceptable, or should \texttt{dt} be a
  hybrid logical clock that never runs behind anything its shard has seen?
\item \textbf{Space.} Nothing is ever lost, so superseded and deleted versions keep
  their blobs. Is anything ever reclaimed, and under whose policy?
\item \textbf{Volume.} Many cameras, all day: a record is visible only once its
  blob is here, so every shard that replicates a camera folder holds all of its
  video. Should a folder's policy be able to keep its raw video at its origin
  while the derived text replicates, with the video fetched by range on demand?
  That relaxes ``a visible record has its blob'' for those folders.
\item \textbf{A new shard.} Its cursors start at zero, and it pulls every peer's
  whole feed; the waiting bound of 100000 records throttles that. Should a new
  shard start from a copy instead?
\item \textbf{A seat's home shard.} Does a seat write only on its home shard,
  which would make concurrent edits of its files rare, or anywhere?
\item \textbf{Retried uploads.} How long is \texttt{request\_id} remembered, and
  is a retried upload on the other island a second file (as specified) or the
  same one?
\end{enumerate}

\section{Content, jobs and feeds}
\begin{enumerate}
\item \textbf{Uploading a folder.} With \texttt{curl}, one file at a time in
  parts, or a tar that the shard unpacks, each entry passing the same checks?
\item \textbf{The job chain.} Which jobs run on which kinds of file, and may a
  derived file trigger further jobs (\texttt{x.mp4--audio.wav--lyrics.vtt})?
  Jobs run only where the blob was created, so the chain runs once.
\item \textbf{Where jobs run.} LLM jobs want GPUs, which may be one big machine
  per island. May a worker sit on another machine in the island, and how does it
  reach the blob and return its outputs?
\item \textbf{Which models may see a file?} A folder's policy could supply the
  models, or model origins, that jobs may run on its files (an earlier map system
  had a rule against Chinese models); a job whose model is not allowed does not run. Per
  file, per job, or both?
\item \textbf{New questions of old footage.} After an event, people want new jobs
  run over weeks of video. Can jobs be requested again later, over a range of old
  files, and bounded so they do not starve the live cameras?
\item \textbf{Annotations more sensitive than their source.} An annotation that
  names a suspect gets the video's permissions by default; its own
  \texttt{--policy.rego} can narrow them. Should some jobs write their outputs
  with a narrower policy from the start?
\item \textbf{Cameras never finish.} A file exists only when it is complete, and a
  camera's stream never is. Does a camera upload a segment a minute, each its own
  file, with a wall-clock start in a clock the cameras agree on?
\item \textbf{Moving objects.} A track that reports every second is a new version
  every second. Is a track a file whose versions are its positions, or a series of
  segment files? Either way the map needs the latest position of thousands of
  objects at once.
\item \textbf{Feeds that cost money.} Every paid external feed needs a budget it
  cannot exceed: a cap on requests or spend per period, after which it stops and
  says so. Who sets it, per feed or per island?
\end{enumerate}

\section{Search and speed}
\begin{enumerate}
\item \textbf{Search is a service.} Indexed things live in a database behind an
  API, so search is mounted under \texttt{/apis/}, filled by subscribing to the
  shard's feed. A hit's name needs \texttt{list}; its snippet and position (seconds
  into a video) need \texttt{read}; a count of hidden hits must not show. Which
  engine: Elasticsearch, as the original should have used, or something embedded?
\item \textbf{How fast is a query?} A listing, a feed page, a search: what is the
  target for the first result, under what load? An earlier map system's ten
  seconds is the number to beat by a wide margin.
\end{enumerate}

\section{Subscribers}
A UI subscribes to a shard's feed exactly as a peer does. There is one feed: the
records readable on that shard, in the order they became readable, read from a
cursor in that shard's clock. A peer is the subscriber whose token can read
everything; a browser is one whose token reads less. Replication is invisible to
the UI, because an import is one more event in the feed. Two things follow for
free:
\begin{itemize}\tightlist
\item \textbf{Failover is a second cursor.} A browser keeps a cursor per shard,
  \texttt{\{a: 50, b: 99\}}, exactly as a shard keeps one per peer. Moving to the
  other island resumes from that island's cursor, and a record seen on both is
  dropped by its \texttt{(id, dt)}.
\item \textbf{Order is the same function.} Events arrive in commit order, so a
  concurrent version can arrive after a newer one. The UI keeps the greatest
  \texttt{(dt, shard)} per \texttt{id}, the same function every shard computes.
\end{itemize}
What comes only from a token that cannot read everything (a policy change
reaching back in time, and a cursor that reveals hidden activity) is in the
security questions above.

\section{Agreed, not yet in Lean}
\begin{itemize}\tightlist
\item \textbf{The tree and its policies.} Content at ordinary URLs under
  \texttt{/files/<org>/<seat>}; metadata built as a table from the root down, each
  folder's \texttt{.--policy.rego} and then the file's
  \texttt{<name>--policy.rego} overwriting the entries they set; a seat's home
  folder and its policy made on the first verified token. Writing into an empty
  folder may supply its \texttt{.--policy.rego} first; otherwise the folder gets
  its parent's. Writing a file creates the missing folders on its path, like
  \texttt{mkdir -p}, if the walk to the nearest existing folder allows it. The
  specification still uses \texttt{plans.pdf.rego} and
  \texttt{/files/\{id\}/\{dt\}/content}.
\item \textbf{\texttt{list} beside \texttt{read} and \texttt{write}.} 404 without
  \texttt{list}, 403 with \texttt{list} but not \texttt{read}; \texttt{list}
  narrows down the path and \texttt{read} implies it. The \texttt{Decision}
  message in \texttt{Requirements.proto} has no \texttt{list} yet, and
  \texttt{AccessDenied} in \texttt{Session.lean} (\texttt{Download},
  \texttt{Forward}, \texttt{Socket}) allows only 401 and 404, not 403.
\item \textbf{Names are content.} A name the caller may not list never reaches
  them, and every channel that could carry one needs a negative requirement and a
  test:
  \begin{itemize}\tightlist
  \item listings, feed pages, history and search results;
  \item a 404 body or header that echoes the requested path;
  \item server and access logs, which must record ids rather than paths, or be
    protected at the level of what they record;
  \item the \texttt{Referer} header a viewer frame sends, and browser history:
    content pages are served with \texttt{Referrer-Policy: no-referrer};
  \item metadata and job outputs, whose names embed their source's name: by
    default they share its permissions, and a policy may not loosen that;
  \item the paths of folders between the root and a readable file.
  \end{itemize}
  A caller who already knows a hidden path learns nothing from asking for it: the
  answer is the same as for a path that never existed.
\item \textbf{No classification system in the source.} The server never names a
  marking, a rank or a rule; they all come from policy files. Tests: the same
  binary serves a NATO-marked tree and a GDPR tree, each with the right banners
  and decisions, and the source contains no marking strings.
\item \textbf{Trust is chosen per shard.} Each shard picks its peers and its token
  issuers. The properties ``Converge after repair'' and ``Same view everywhere
  after replication'' in \texttt{Requirements.lean} still speak of every shard;
  they must be restated per token: any two shards that accept the token, and hold
  the same records, give it the same answers; and the shards that accept a token
  converge with each other after repair.
\item \textbf{The door.} \texttt{ingressHttps} serves \texttt{/files/} and forwards
  \texttt{/apis/} mounts to internal services. Its protocols, \texttt{Forward} and
  \texttt{Socket}, are in \texttt{Session.lean} and checked; its negative
  requirements are not yet: no service is reachable except through it; no request
  reaches a service without a verified token and an allowing walk; no identity
  header from a client reaches a service; every request is counted per mount.
\item \textbf{Apps read, APIs write.} An uploaded app is read-only to its users.
  Every API that writes identifies the user and checks their permissions on every
  call, with negative requirements and tests like the core: no call without a
  verified identity; no write the policy at its path does not allow; no input
  reaching a query, a shell or a path unescaped; no answer that names what the
  caller may not list; every write recorded with who made it.
\item \textbf{Two halves, one pattern.} \texttt{/files/} for data and apps,
  \texttt{/apis/} for services users mount. An API serves its data as a feed read
  by cursors, with no state per user, like the record feed.
\item \textbf{One feed for browsers and peers.} \texttt{FeedRequest} and
  \texttt{RecordsRequest} become one message: records past a cursor, in the
  shard's \texttt{visible\_dt} order, filtered by the caller's token. The two
  tasks, and their properties, are still separate in \texttt{Requirements.lean}.
\item \textbf{Content brings its viewer.} The shard serves uploaded HTML, scripts
  and styles as a web site, so a folder can be an app; deep links such as
  \texttt{?ts=40} are the viewer's own convention.
\item \textbf{Derived files are named \texttt{source--x}.} With two dashes they are
  metadata about their source and share its permissions by the naming convention.
  \texttt{Requirements.lean} still names them with one dash.
\item \textbf{Every derived thing answers ``what is this?''.} A report entry, a
  map object or a search hit links back, through its name, to the raw file it
  came from, which the caller may open only if they may read it.
\item \textbf{Replay costs no blob.} A resent record already held, by its
  \texttt{(id, dt)}, causes no \texttt{BlobRequest}. It is in
  \texttt{Implementation.lean}, but not yet a requirement a test can check.
\item \textbf{Seeking is fast everywhere.} The first byte of any range arrives
  within a fixed time, whatever its offset in the file: the 9000 MB video that
  seeked instantly. It is not yet a performance requirement.
\item \textbf{Simulated sources for stress tests.} Like the fake airport that made
  planes fly around an earlier map: simulated cameras, feeds and tracks, at
  rates set by the test, drive the performance requirements (uploads per second,
  the latest positions of thousands of tracks, the first byte of a search) on both
  islands, through a partition and its repair.
\end{itemize}

\section{Limits of the tools}
\begin{itemize}\tightlist
\item The models in \texttt{Islands.lean} are written by hand; they are not yet
  derived from the session types.
\item The MCP server refuses arguments over 64 KiB (\texttt{maxArgumentBytes} in
  LeanFM's \texttt{Logic.lean}; HTTP itself allows 8 MiB) and models over 4096
  nodes. Every model here is scaled down to fit: two shards, two or three
  records, two lanes. \texttt{repl-cursor} is 61 KB. A larger bound with a time
  budget would let the models grow.
\item A nested \texttt{EF} witness stops at its first conjunct.
  \texttt{mcp\_islands.py} joins the rest by asking again from where it stopped.
\item Liveness holds only under fair delivery. The server shows the lassos where
  it fails, and the requirements state the assumption instead of hiding it.
\end{itemize}

\appendix
\chapter{The project Makefile}
The project needs this file at \texttt{projects/IslandUploads/Makefile}, and
\texttt{book/Makefile} is a copy of PaxosKV's. Copying from a PDF turns tabs into
spaces, so recipe lines are written with \texttt{>} and the last line turns
them into tabs. From the LeanFM root:
\begin{Verbatim}[fontsize=\small]
cp projects/PaxosKV/book/Makefile projects/IslandUploads/book/Makefile
cat > projects/IslandUploads/Makefile <<'MAKEFILE'
LEANFM := $(abspath ../..)
.PHONY: all check spec-check session models mcp book clean
all: check book
spec-check:
> cd $(LEANFM) && python3 scripts/chatui_specification.py $(CURDIR) \
    Project.Requirements.requirements Project.Implementation.implementation
session:
> cd $(LEANFM) && lake build LeanFM.Session && \
    lake env lean $(CURDIR)/spec/Session.lean
models:
> cd $(LEANFM) && lake env lean --run $(CURDIR)/spec/Islands.lean
mcp: models
> python3 scripts/mcp_islands.py
> python3 scripts/mcp_report.py
check: spec-check session
book:
> $(MAKE) -C book
clean:
> $(MAKE) -C book clean
MAKEFILE
perl -pi -e 's/^> /\t/' projects/IslandUploads/Makefile
\end{Verbatim}
